\scanpage{104}
\mspage{51}
qui sont induits par les morphismes de normalisation entre m.graphes associés aux graphes monoponctuels. Bien entendu, pour pouvoir reconstituer bel et bien les \inserted{effets \deleted{\unsure{sur}} en $\Pi_1M$ des} morphismes de normalisation, et leurs propriétés de commutation aux iso., il faut utiliser les \deleted{\unsure{covariances}} \unsure{commutation|covariances} de (148) à
\[
\unsure{\mathfrak S|\mathcal S}_{I,I'}\simeq\unsure{\mathfrak S|\mathcal S}_{I'}\times\unsure{\mathfrak S|\mathcal S}_{I\setminus I'}\qquad\ldots
\]

\uline{Programme} Déterminer (à équivalence ``canonique'' près) tous les groupoïdes $\Pi_1M_{g,I}$, avec leur dépendance fonctorielle pour \inserted{$g$ fixe,} $I$ variable (des mono. $I'\hookrightarrow I$) via les morphismes de normalisation (148).
\margin{par ``\unsure{gén.|génér.} et relations'' — \unsure{définies|définis} en termes purement algébriques\ldots}

De plus, pour tt MD-graphe $G$, déterminer l'extension $S\Pi_1M_G$ de
\[
\Pi_1M_G\simeq\prod_{\alpha\in S(G)}\Pi_1M_{\unsure{\alpha|g_\alpha},\widehat I_\alpha}\quad\text{par }\mathbb Z^{\widetilde A(G)},
\]
et pour un morphisme de générisation (p.ex. un iso!) \deleted{donné} $G\longrightarrow G'$ de MD-graphes, \deleted{définir}\inserted{déterminer} le morphisme de groupoïdes de générisation correspondant
\[
S\Pi_1M_G\longrightarrow S\Pi_1M_{G'}.
\]
\margin{\figure{\begin{tikzpicture}\draw (0,0)--(.05,-8);\end{tikzpicture}}}
\margin{103}
